Dues esferes conductores idèntiques i suficientment petites per a ser considerades puntuals estan unides per una molla. La constant elàstica de la molla és $10\ \mathrm{N/m}$. El conjunt es col·loca sobre una taula que és elèctricament aïllant. A més, no hi ha fricció entre la taula i el conjunt de les dues esferes i la molla. Les esferes estan separades per una distància de $0,40\ \mathrm{m}$ quan no estan carregades (la molla no fa cap força). Carreguem amb la mateixa càrrega positiva les dues esferes amb un generador fins que la distància entre elles sigui $1,00\ \mathrm{m}$.

\figura[0.35]{fig1.png}

\begin{apartat}{1,25}
Quan les esferes estan carregades, quina és la força aplicada per la molla sobre cadascuna de les esferes? Quina és la càrrega de cadascuna de les esferes?
\end{apartat}

\begin{apartat}{1,25}
Col·loquem un electró a $1\ \mathrm{m}$ de cadascuna de les dues esferes (equidistant a les dues esferes, tal com indica la figura). Calculeu el mòdul del camp elèctric que actua sobre l'electró i el mòdul de l'acceleració en aquest instant. Sobre la figura, representeu la direcció i sentit del camp elèctric i de l'acceleració en el punt on es troba l'electró.
\end{apartat}

\dades{%
  $k_\mathrm{e} = \dfrac{1}{4\pi\varepsilon_0} = 8,99\times 10^{9}\ \mathrm{N\,m^2\,C^{-2}}$.\\[2pt]
  $m_\mathrm{e} = 9,11\times 10^{-31}\ \mathrm{kg}$.\\
  $|e| = 1,602\times 10^{-19}\ \mathrm{C}$.}
